% chapters/09_development.tex
\chapter{研发过程与工程收尾}

\section{研发原则}

项目采用“先保证安全语义，再接入真实块路径，最后固化复现入口”的顺序：

\begin{enumerate}
    \item 在用户态用 \texttt{BlockDevice} 抽象验证 AES-XTS、Merkle tree 与错误边界；
    \item 通过 \texttt{Virtio\allowbreak Blk\allowbreak Adapter} 接入 Asterinas BIO，不绕过真实下层设备；
    \item 分别完成 crypt、verity 与组合栈，再接入受保护 rootfs；
    \item 使用 Linux 兼容性、KVM 正反例和物理损坏测试作为合入门禁；
    \item 性能数据必须通过 target、参数、parser 与完成语义审计。
\end{enumerate}

\section{关键工程问题}

\subsection{块大小与扇区语义}

上层文件系统通常以 4 KiB 块发起 I/O，而 dm-crypt plain64 IV 以 512 字节扇区号
为单位。实现将每个块拆分为连续扇区，并同时考虑 IV offset 与 data offset；所有
加减法都进行越界检查，避免整数溢出改变物理位置。

\subsection{BIO 生命周期}

mapper 只在 lower BIO 完成后完成 upper BIO。读路径先获得密文，再按扇区解密并写回
上层 segment；写路径复制明文、加密临时缓冲并提交 lower BIO。错误状态原样传播，
不会把失败转换为成功，也不会 fallback 到未保护设备。

\subsection{verity 信任边界}

root hash 由启动参数提供，镜像内数据和 hash tree 均视为不可信。认证后 hash cache
只有在完整路径验证成功后才插入，数据块、叶子 hash block 或中间 hash block 的单字节
变化都必须在交付前被拒绝。

\section{Asterinas 补丁演进}

\begin{table}[H]
    \centering
    \caption{内核移植补丁职责}
    \begin{tabular}{p{0.15\textwidth}p{0.72\textwidth}}
        \toprule
        \textbf{补丁} & \textbf{职责} \\
        \midrule
        0001 & Device Mapper 组件骨架、crypt/verity target 与 VirtIO 适配 \\
        0002 & 启动参数、mapper 注册和 ext2 魔数自检 \\
        0003 & initramfs 挂载与 \texttt{pivot\_\allowbreak root} \\
        0004 & cryptsetup 控制面与 mapper 生命周期 \\
        0005 & 正式 benchmark 参数和内核统计接口 \\
        0006 & 性能阶段计数与可复核元数据 \\
        0007 & \texttt{dm-crypt} $\to$ \texttt{dm-verity} Combined 组合栈 \\
        \bottomrule
    \end{tabular}
\end{table}

当前仓库 gitlink 已包含正式集成版本；\texttt{asterinas\_\allowbreak port/\allowbreak patches} 用于说明移植
边界与复核历史，不要求评委在正常 checkout 后重复应用。

\section{渐进式保护路径}

\begin{table}[H]
    \centering
    \caption{五阶段验证路径}
    \begin{tabular}{p{0.10\textwidth}p{0.20\textwidth}p{0.58\textwidth}}
        \toprule
        \textbf{阶段} & \textbf{名称} & \textbf{数据流} \\
        \midrule
        1 & Plain & virtio-blk $\to$ ext2 \\
        2 & Passthrough & virtio-blk $\to$ Passthrough\allowbreak Device $\to$ ext2 \\
        3 & dm-verity & virtio-blk $\to$ DmVerity\allowbreak Device $\to$ ext2 \\
        4 & dm-crypt & virtio-blk $\to$ DmCrypt\allowbreak Device $\to$ ext2 \\
        5 & Combined & virtio-blk $\to$ DmCrypt\allowbreak Device $\to$ DmVerity\allowbreak Device $\to$ ext2 \\
        \bottomrule
    \end{tabular}
\end{table}

\begin{sloppypar}
Combined 的镜像构建与启动入口为 \texttt{scripts/\allowbreak build\_\allowbreak stack\_\allowbreak rootfs.sh} 和
\texttt{scripts/\allowbreak run\_\allowbreak dm\_\allowbreak stack.sh}。用户态 \texttt{dmctl demo-stack} 使用相同装饰器顺序。
\end{sloppypar}

\section{测试代码收口}

最终版本只保留当前正式工作流使用的构建、正确性和性能 runner。一次性调试脚本、旧
跨 guest 比较器、失败重试和早期小矩阵不进入最终目录。发布检查执行以下门禁：

\begin{lstlisting}[basicstyle=\ttfamily\small, breaklines=true]
cargo test --workspace
python3 scripts/check_benchmark_package.py
python3 scripts/verify_kvm_evidence.py
python3 scripts/review_release_check.py
\end{lstlisting}

静态审查要求 Python 不存在未使用函数/import，shell 函数必须存在实际调用；生成镜像
和缓存位于 \texttt{target}，不上传到代码仓库。

\section{已知边界}

\begin{itemize}
    \item 当前主路径聚焦 AES-256-XTS、plain64 与 SHA-256，不宣称覆盖 Linux 的全部 target 选项；
    \item 用户态 verity tree 用于快速测试，Asterinas 内核路径使用 on-disk hash tree；
    \item dm-verity 为只读目标，写入 workload 明确记为 N/A；
    \item 不同平台的性能环境分别报告，避免把宿主机差异解释为产品收益；
    \item 密钥注入的生产化方案仍需结合可信启动、远程证明和密钥管理系统。
\end{itemize}

\section{小结}

研发过程最终收敛为单一正式代码路径、单一性能矩阵生成器和最小可复核证据链。核心安全
实现没有为了改善跑分而放宽验证条件，历史无效数据也不会进入最终公开结论。